\section{Thảo luận}
\label{sec:discussion}

Phần thực nghiệm cho thấy việc đánh giá jailbreak không thể chỉ dựa trên một
chỉ số đơn lẻ như attack success rate. Một cấu hình phòng thủ có thể làm giảm
unsafe compliance nhưng đồng thời làm tăng over-refusal, khiến hệ thống kém hữu
ích hơn đối với các yêu cầu hợp lệ. Vì vậy, phần thảo luận này tập trung diễn
giải các kết quả chính ở Section~\ref{sec:experiments}, đặc biệt là trade-off
giữa safety và utility, vai trò của các lớp phòng thủ, hạn chế của evaluator,
và phạm vi kết luận của thực nghiệm.

\subsection{Ý nghĩa của kết quả thực nghiệm}

Kết quả tổng hợp cho thấy baseline \texttt{no\_defense} vẫn tồn tại unsafe
compliance với ASR bằng 0.133. Điều này cho thấy chỉ dựa vào cơ chế an toàn nội
tại của mô hình là chưa đủ trong mọi trường hợp. Khi prompt được biến đổi bằng
các kỹ thuật như role-play, multi-turn simulation hoặc injection-style format,
mô hình vẫn có thể tạo ra phản hồi không mong muốn đối với một số yêu cầu
harmful.

Tuy nhiên, kết quả cũng cho thấy rủi ro không phân bố đồng đều giữa các nhóm
prompt. Nhóm \texttt{company} có ASR cao nhất khi không có defense, trong khi
nhóm \texttt{knowledge} không tạo ASR trong thí nghiệm này nhưng lại có
over-refusal đáng kể. Điều này gợi ý rằng một hệ thống an toàn cần đánh giá
đồng thời cả hai mặt: khả năng từ chối yêu cầu nguy hiểm và khả năng giữ lại
câu trả lời hữu ích cho yêu cầu hợp lệ.

Một điểm quan trọng là attack success không phải dạng lỗi duy nhất. Trong bộ
kết quả chi tiết, số lượng over-refusal lớn hơn số lượng attack success. Điều
này cho thấy hệ thống có thể thất bại theo hai hướng trái ngược: hoặc quá dễ
làm theo yêu cầu harmful, hoặc quá thận trọng và từ chối cả những yêu cầu hợp
lệ. Cả hai dạng lỗi đều ảnh hưởng đến chất lượng của một hệ thống LLM thực tế.

\subsection{Safety--utility trade-off}

Kết quả thực nghiệm minh họa rõ trade-off giữa safety và utility. Cấu hình
\texttt{output\_moderation} đưa ASR về 0.000 trong khi vẫn giữ Utility Rate ở
mức 0.613. Ngược lại, \texttt{input\_filter} và
\texttt{defense\_in\_depth} cũng giúp giảm ASR nhưng làm ORR tăng lên 0.573 và
Utility Rate giảm còn 0.427.

Điều này cho thấy không phải mọi defense làm giảm ASR đều là lựa chọn tốt. Nếu
một lớp phòng thủ chặn quá nhiều prompt hợp lệ, hệ thống có thể an toàn hơn
theo nghĩa hẹp nhưng kém hữu ích hơn trong thực tế. Trong các ứng dụng như trợ
lý học thuật, phân tích tài liệu hoặc hỗ trợ lập trình, over-refusal cao có thể
làm người dùng không nhận được câu trả lời cho các câu hỏi hợp pháp.

Trong demo này, \texttt{output\_moderation} đạt trade-off tốt nhất vì nó kiểm
tra phản hồi sau khi mô hình đã sinh output. Cách tiếp cận này ít phụ thuộc hơn
vào từ khóa xuất hiện trong input, do đó giảm nguy cơ chặn nhầm các prompt học
thuật hoặc prompt an toàn có chứa thuật ngữ nhạy cảm. Tuy nhiên, kết quả này
không có nghĩa output moderation luôn tốt hơn input filtering trong mọi hệ
thống. Trong môi trường production, input filtering vẫn cần thiết để ngăn một
số yêu cầu rõ ràng không phù hợp đi vào mô hình, đặc biệt khi chi phí gọi model
cao hoặc khi cần giảm rủi ro ngay từ đầu.

\subsection{Vai trò của input filtering và output moderation}

Kết quả cho thấy input filtering có lợi thế là đơn giản, dễ triển khai và có
thể chặn sớm một số prompt có dấu hiệu rủi ro. Tuy nhiên, rule-based input
filtering trong demo này gây over-refusal cao. Nguyên nhân là nhiều prompt hợp
lệ về AI safety, jailbreak hoặc an toàn hóa chất vẫn chứa các từ khóa nhạy cảm.
Nếu filter chỉ dựa trên keyword, hệ thống khó phân biệt giữa thảo luận học thuật
và yêu cầu gây hại thực sự.

Ngược lại, output moderation đánh giá phản hồi sau khi model đã sinh nội dung.
Điều này cho phép hệ thống chỉ chặn khi output thực sự có dấu hiệu unsafe
compliance. Trong thực nghiệm, output moderation giảm ASR về 0.000 mà không làm
giảm Utility Rate so với baseline. Đây là kết quả quan trọng vì nó cho thấy lớp
kiểm soát sau sinh có thể giữ lại nhiều câu trả lời hợp lệ hơn.

Tuy nhiên, output moderation cũng có hạn chế. Nếu moderation model hoặc rule
không đủ mạnh, nội dung không an toàn có thể lọt qua. Ngoài ra, output
moderation xảy ra sau khi model đã sinh phản hồi, nên trong một số hệ thống có
tool-use hoặc streaming output, việc chặn sau sinh có thể không đủ an toàn.
Do đó, trong các hệ thống thực tế, defense-in-depth vẫn cần thiết, nhưng các
lớp defense phải được thiết kế tinh tế hơn so với keyword filtering đơn giản.

\subsection{Vấn đề over-refusal}

Over-refusal là một trong những hiện tượng nổi bật trong thực nghiệm. Nhóm
\texttt{hazardous} có ORR tăng mạnh khi áp dụng input filtering, trong khi nhóm
\texttt{knowledge} cũng có ORR đáng kể ngay cả khi không có defense. Điều này
cho thấy các câu hỏi hợp lệ nhưng chứa thuật ngữ nhạy cảm có thể bị xử lý quá
thận trọng.

Trong bối cảnh AI safety, over-refusal đặc biệt đáng chú ý vì nhiều câu hỏi học
thuật về jailbreak, prompt injection, guardrail hoặc benchmark có thể dùng cùng
từ vựng với các prompt nguy hiểm. Nếu hệ thống không hiểu ngữ cảnh, nó có thể
từ chối cả các yêu cầu phục vụ nghiên cứu, giáo dục hoặc đánh giá an toàn. Điều
này làm giảm khả năng hỗ trợ người dùng và có thể khiến hệ thống kém hữu ích
trong các tác vụ phân tích.

Kết quả này củng cố nhận định rằng đánh giá LLM safety cần luôn đi kèm đánh giá
utility. Một mô hình chỉ từ chối nhiều hơn không nhất thiết là một mô hình tốt
hơn. Hệ thống cần học cách phân biệt giữa yêu cầu có mục đích gây hại và yêu
cầu thảo luận an toàn ở mức khái niệm.

\subsection{Độ tin cậy của evaluator}

Evaluator trong demo hiện tại là rule-based. Cách tiếp cận này có ưu điểm là
đơn giản, dễ kiểm tra và dễ tái lập. Tuy nhiên, nó cũng là một trong những giới
hạn lớn nhất của thực nghiệm. Một phản hồi dài có thể bắt đầu bằng câu từ chối
nhưng vẫn chứa phần giải thích không phù hợp phía sau. Ngược lại, một phản hồi
có chứa từ khóa nhạy cảm không nhất thiết là unsafe nếu nó đang giải thích ở
mức khái quát hoặc từ chối một cách an toàn.

Do đó, nhãn \texttt{model\_label} và \texttt{final\_label} trong experiment nên
được xem là nhãn gần đúng. Chúng đủ để minh họa pipeline và tạo kết quả định
lượng ban đầu, nhưng chưa thay thế được human evaluation hoặc LLM-as-a-judge.
Trong một benchmark nghiêm ngặt hơn, cần có rubric gán nhãn rõ ràng, nhiều
người đánh giá độc lập, hoặc một judge model được kiểm định.

Vấn đề này cũng liên quan đến manual web evaluation. Các cột \texttt{Passed\_*}
trong file manual được giữ lại như nhãn nguồn, nhưng không được tự động diễn
giải thành \textit{unsafe compliance} hay \textit{safe refusal}. Vì vậy,
\texttt{passed\_1\_rate} không được xem là ASR. Manual evaluation chỉ đóng vai
trò bổ sung, giúp xác định những case nên đọc thủ công sâu hơn.

\subsection{Hạn chế do tài nguyên}

Do giới hạn về tài nguyên tính toán, quota API và thời gian thực hiện, thí
nghiệm có phạm vi tương đối nhỏ. Bộ prompt chỉ gồm 30 prompt gốc, không thể đại
diện cho toàn bộ không gian jailbreak prompt. Target model chính được gọi qua
OpenRouter API, nên kết quả phụ thuộc vào model, thời điểm gọi API, chính sách
của provider và cấu hình generation.

Ngoài ra, nhóm không chạy lại các official implementation nặng hơn như GCG,
PAIR hoặc AutoDAN. Các phương pháp này thường cần nhiều lần gọi model, truy cập
chi tiết hơn vào mô hình hoặc tài nguyên GPU. Vì vậy, thực nghiệm trong báo cáo
này được định vị là demo tự xây dựng theo Cách 2, không phải reproduction đầy
đủ của một paper cụ thể.

Các defense cũng được cài đặt ở mức đơn giản. Input filter dựa trên keyword,
output moderation dựa trên rule hoặc nhãn suy luận, và defense-in-depth là kết
hợp của hai thành phần trên. Điều này đủ để minh họa trade-off nhưng chưa phản
ánh đầy đủ hệ thống guardrail production, nơi thường sử dụng classifier chuyên
dụng, policy engine, audit log và human review.

\subsection{Hàm ý cho thiết kế hệ thống thực tế}

Từ kết quả thực nghiệm, có thể rút ra một số hàm ý cho thiết kế hệ thống LLM
thực tế.

Thứ nhất, defense nên được đánh giá bằng nhiều chỉ số. Nếu chỉ tối ưu ASR, hệ
thống có thể trở nên quá nghiêm ngặt và làm giảm utility. Ngược lại, nếu chỉ
tối ưu utility, hệ thống có thể bỏ qua các trường hợp unsafe compliance. Một
đánh giá cân bằng cần bao gồm ít nhất ASR, ORR, Utility Rate và phân tích lỗi.

Thứ hai, defense-in-depth là cần thiết nhưng không nên được hiểu là cộng nhiều
rule đơn giản vào nhau. Nếu mỗi lớp đều chặn rộng, tổng thể hệ thống sẽ có
over-refusal cao. Defense-in-depth hiệu quả cần phân tầng theo ngữ cảnh: input
guard để chặn yêu cầu rõ ràng nguy hiểm, output guard để kiểm tra nội dung
được sinh ra, và cơ chế escalation cho các trường hợp không chắc chắn.

Thứ ba, các nhóm prompt khác nhau cần cách xử lý khác nhau. Prompt liên quan
đến thông tin nội bộ công ty có rủi ro khác với prompt học thuật về jailbreak
hoặc prompt an toàn hóa chất. Một policy chung dựa trên keyword có thể không đủ
tinh tế để xử lý mọi nhóm nội dung.

\subsection{Hướng cải thiện}

Nếu có thêm tài nguyên, thực nghiệm có thể được mở rộng theo các hướng sau:

\begin{itemize}[leftmargin=*]
    \item \textbf{Mở rộng dataset}: tăng số lượng prompt, bổ sung benchmark
    công khai và phân tầng rõ hơn theo mức độ rủi ro.

    \item \textbf{So sánh nhiều target models}: chạy cùng pipeline trên nhiều
    mô hình khác nhau để đánh giá độ ổn định của kết quả.

    \item \textbf{Cải thiện evaluator}: thay rule-based evaluator bằng
    LLM-as-a-judge, classifier chuyên dụng hoặc human evaluation có rubric.

    \item \textbf{Tích hợp defense mạnh hơn}: thử nghiệm các guard model,
    contextual moderation hoặc các defense được đề cập trong tutorial như
    smoothing, adversarial training hoặc circuit breaker.

    \item \textbf{Phân tích sâu error cases}: đọc thủ công các case unsafe
    compliance và over-refusal để xác định nguyên nhân cụ thể.

    \item \textbf{Mở rộng sang agentic systems}: kiểm thử prompt injection
    trong bối cảnh tool-use, memory, retrieval hoặc multi-agent workflow.
\end{itemize}

Nhìn chung, phần thảo luận cho thấy kết quả thực nghiệm không nên được hiểu như
một kết luận tổng quát về mọi mô hình hoặc mọi phương pháp jailbreak. Giá trị
chính của demo là minh họa được cách xây dựng pipeline đánh giá, cách đo lường
trade-off giữa safety và utility, và các vấn đề thực tế phát sinh khi triển khai
defense trong hệ thống LLM.
